\chapter{Chapter05 — Prescribed Fourier Support}
%%% generated by the blueprint pipeline from FABL.Chapter05.PrescribedFourierSupport
\section{Overview}

Declarations of the Lean module \texttt{FABL.Chapter05.PrescribedFourierSupport}.

\section{Declarations}

\begin{definition}[Singleton-indicator Fourier projection]
\lean{FABL.singletonIndicatorFourierProjection}
\label{FABL.singletonIndicatorFourierProjection}
\leanok
\uses{FABL.fourierCoeff, FABL.indicatorPolynomial, FABL.monomial}
Let $\mathcal F$ be a finite family of subsets of $\{0,\ldots,n-1\}$ and let
$a\in\{−1,1\}^n$. The singleton-indicator Fourier projection is the function sending
$x\in\{−1,1\}^n$ to $\sum_{S\in\mathcal F}\widehat{\mathbf 1_{\{a\}}}(S)\,\chi_S(x)$,
where $\widehat{\mathbf 1_{\{a\}}}(S)$ is the Fourier coefficient at $S$ of the
indicator of the singleton $\{a\}$ and $\chi_S(x)=\prod_{i\in S}x_i$.
\end{definition}

\begin{definition}[Prescribed Fourier kernel]
\lean{FABL.prescribedFourierKernel}
\statementsource{ghsz-combinatorial-list-decoding}{pdf-6bd33ae73ade-p007-b009}
\label{FABL.prescribedFourierKernel}
\leanok
\uses{FABL.singletonIndicatorFourierProjection}
For a finite family $\mathcal F$ of subsets of $\{1,\ldots,n\}$ and a point
$a\in\{-1,1\}^n$, define a function on the Boolean cube by $K_{\mathcal
F,a}(x)=\dfrac{2^n}{|\mathcal F|}\,P_{\mathcal F}(a,x)$, where $P_{\mathcal F}$ is the
singleton-indicator Fourier projection associated with $\mathcal F$.
\end{definition}

\begin{theorem}[Fourier coefficient of an indicator polynomial]
\lean{FABL.fourierCoeff_indicatorPolynomial}
\label{FABL.fourierCoeff_indicatorPolynomial}
\leanok
\uses{FABL.fourierCoeff, FABL.indicatorPolynomial, FABL.indicatorPolynomial_eq_ite, FABL.monomial}
\difficulty{3}
Let $n$ be a nonnegative integer and let $a\in\{\pm1\}^n$. For every set of coordinates
$S\subseteq\{0,\ldots,n-1\}$, the Fourier coefficient at $S$ of the polynomial
$x\mapsto\prod_{j=0}^{n-1}(1+a_jx_j)/2$ equals $2^{-n}$ times the monomial $\prod_{j\in
S}a_j$.
\end{theorem}

\begin{theorem}[Prescribed Fourier kernel formula]
\lean{FABL.prescribedFourierKernel_apply}
\label{FABL.prescribedFourierKernel_apply}
\leanok
\uses{FABL.fourierCoeff_indicatorPolynomial, FABL.monomial, FABL.prescribedFourierKernel, FABL.singletonIndicatorFourierProjection}
\difficulty{4}
Let $n$ be a natural number, let $\mathcal F$ be a nonempty finite family of subsets of
$\{0,\ldots,n-1\}$, and let $a,x\in\{−1,1\}^n$. The prescribed Fourier kernel at $a$ and
$x$ is $1/|\mathcal F|$ times the sum, over $S\in\mathcal F$, of the monomial indexed by
$S$ evaluated at $a$ multiplied by that monomial evaluated at $x$.
\end{theorem}

\begin{theorem}[Prescribed Fourier kernel on the diagonal]
\lean{FABL.prescribedFourierKernel_apply_self}
\label{FABL.prescribedFourierKernel_apply_self}
\leanok
\uses{FABL.monomial_sq, FABL.prescribedFourierKernel, FABL.prescribedFourierKernel_apply}
\difficulty{3}
Let $n$ be a natural number, let $\mathcal F$ be a nonempty finite family of subsets of
an $n$-element coordinate set, and let $a\in\{−1,1\}^n$. The prescribed Fourier kernel
associated with $\mathcal F$, evaluated at $(a,a)$, equals $1$.
\end{theorem}

\begin{theorem}[Fourier coeff prescribed fourier kernel]
\lean{FABL.fourierCoeff_prescribedFourierKernel}
\statementsource{ghsz-combinatorial-list-decoding}{pdf-6bd33ae73ade-p007-b009}
\label{FABL.fourierCoeff_prescribedFourierKernel}
\leanok
\uses{FABL.expect_monomial_mul, FABL.fourierCoeff, FABL.monomial, FABL.prescribedFourierKernel, FABL.prescribedFourierKernel_apply}
\difficulty{5}
% dataset-pipeline: stub -- prose pending build_informal_statements.py
Statement pending regeneration.
\end{theorem}

\begin{theorem}[Second moment of a prescribed Fourier kernel]
\lean{FABL.expect_sq_prescribedFourierKernel}
\label{FABL.expect_sq_prescribedFourierKernel}
\leanok
\uses{FABL.fourierCoeff, FABL.fourierCoeff_prescribedFourierKernel, FABL.monomial_sq, FABL.parseval, FABL.prescribedFourierKernel, FABL.uniformInner}
\difficulty{5}
Let $n$ be a natural number, let $\mathcal{F}$ be a nonempty family of subsets of
$\{1,\ldots,n\}$, and let $a\in\{-1,1\}^n$. The uniform expectation over
$x\in\{-1,1\}^n$ of the square of the prescribed Fourier kernel associated with
$\mathcal{F}$ and $a$ is $1/|\mathcal{F}|$.
\end{theorem}

\begin{theorem}[Symmetry of the prescribed Fourier kernel]
\lean{FABL.prescribedFourierKernel_comm}
\label{FABL.prescribedFourierKernel_comm}
\leanok
\uses{FABL.prescribedFourierKernel, FABL.prescribedFourierKernel_apply}
\difficulty{3}
Let $n$ be a natural number, let $\mathcal F$ be a nonempty finite family of subsets of
$\{0,\ldots,n-1\}$, and let $a,x$ be sign vectors of length $n$. The prescribed Fourier
kernel evaluated at $(a,x)$ equals the kernel evaluated at $(x,a)$.
\end{theorem}

\begin{theorem}[Squared mass of the prescribed Fourier kernel]
\lean{FABL.sum_ne_sq_prescribedFourierKernel}
\label{FABL.sum_ne_sq_prescribedFourierKernel}
\leanok
\uses{FABL.expect_sq_prescribedFourierKernel, FABL.prescribedFourierKernel, FABL.prescribedFourierKernel_apply_self, FABL.prescribedFourierKernel_comm}
\difficulty{4}
Let $n$ be a natural number, let $\mathcal F$ be a nonempty finite family of subsets of
$\{1,\ldots,n\}$, and let $x\in\{-1,1\}^n$. For each $a\in\{-1,1\}^n$, form the function
whose value at $y$ is the sum over $S\in\mathcal F$ of the Fourier coefficient of the
indicator of $a$, on $S$, times the $S$-monomial of $y$, and multiply this function by
$2^n/|\mathcal F|$; then the sum of its squared values at $x$, over all $a\ne x$, is
$2^n/|\mathcal F|-1$.
\end{theorem}

\begin{definition}[Off-center Fourier sum]
\lean{FABL.prescribedFourierOffCenterSum}
\label{FABL.prescribedFourierOffCenterSum}
\leanok
\uses{FABL.BooleanFunction, FABL.prescribedFourierKernel, FABL.signValue}
For a finite family $\mathcal F$ of subsets of $\{0,\ldots,n-1\}$, a point
$x\in\{-1,1\}^n$, and a Boolean function $f:\{-1,1\}^n\to\{\pm1\}$, define the
off-center sum at $x$ to be $\sum_{a\in\{-1,1\}^n,\ a\ne x} f(a)K_{\mathcal F}(a,x)$,
where $K_{\mathcal F}$ is the prescribed Fourier kernel.
\end{definition}

\begin{theorem}[Uniform law on Boolean functions]
\lean{FABL.uniformBooleanFunction_toMeasure_eq_pi}
\label{FABL.uniformBooleanFunction_toMeasure_eq_pi}
\leanok
\uses{FABL.BooleanFunction, FABL.Sign, FABL.uniformPMF}
\difficulty{3}
For every natural number $n$, choose uniformly at random a function from ${pm1}^n$ to
$pm1$. Its probability measure is the product, over all inputs in ${pm1}^n$, of the
uniform probability measure on $pm1$.
\end{theorem}

\begin{theorem}[Mean of a uniform sign]
\lean{FABL.integral_signValue_uniformPMF}
\label{FABL.integral_signValue_uniformPMF}
\leanok
\uses{FABL.Sign, FABL.integral_uniformPMF_eq_expect, FABL.signValue, FABL.uniformPMF}
\difficulty{3}
If $s$ is uniformly distributed over the two integer signs $\{-1,1\}$, then its expected
value is zero.
\end{theorem}

\begin{theorem}[Weighted Boolean sum tail bound]
\lean{FABL.measure_abs_weightedBooleanSum_ge_le}
\label{FABL.measure_abs_weightedBooleanSum_ge_le}
\leanok
\uses{FABL.BooleanFunction, FABL.Sign, FABL.integral_signValue_uniformPMF, FABL.signValue, FABL.signValue_eq_neg_one_or_one, FABL.uniformBooleanFunction_toMeasure_eq_pi, FABL.uniformPMF}
\difficulty{6}
Let $n$ be a natural number, let $w$ assign a real weight to each point of the Boolean
cube $\\{-1,1\\}^n$, and let $t\\ge 0$. Choose a function
$f:\\{-1,1\\}^n\\to\\{\\pm1\\}$ uniformly at random, and write $f(a)$ for its sign
value. The probability that $\\left|\\sum_{a\\in\\{-1,1\\}^n} w(a)f(a)\\right|\\ge t$ is
at most $2\\exp\\left(-\\frac{t^2}{2\\sum_a w(a)^2}\\right)$.
\end{theorem}

\begin{theorem}[Off-center sum as a weighted sum]
\lean{FABL.prescribedFourierOffCenterSum_eq_weightedSum}
\label{FABL.prescribedFourierOffCenterSum_eq_weightedSum}
\leanok
\uses{FABL.BooleanFunction, FABL.prescribedFourierKernel, FABL.prescribedFourierOffCenterSum, FABL.signValue}
\difficulty{4}
Let $n$ be a natural number, let $\mathcal F$ be a finite family of subsets of
$\{1,\ldots,n\}$, let $x$ be a sign vector of length $n$, and let $f$ map sign vectors
of length $n$ to signs. The sum over all sign vectors $a\ne x$ of the prescribed Fourier
kernel at $(a,x)$ multiplied by the real value of $f(a)$ equals the sum over all sign
vectors $a$ of the same product, with the kernel factor set to zero when $a=x$.
\end{theorem}

\begin{theorem}[Off-center kernel squared sum]
\lean{FABL.sum_sq_prescribedFourierOffCenterWeight}
\label{FABL.sum_sq_prescribedFourierOffCenterWeight}
\leanok
\uses{FABL.prescribedFourierKernel, FABL.sum_ne_sq_prescribedFourierKernel}
\difficulty{3}
Let $n$ be a natural number, let $\mathcal F$ be a nonempty family of subsets of the $n$
coordinates, and fix a sign vector $x$. The sum over all sign vectors $a$, excluding
$a=x$, of the squared prescribed Fourier kernel evaluated at $(a,x)$ is $2^n/|\mathcal
F|-1$.
\end{theorem}

\begin{theorem}[Variance bound for a large family]
\lean{FABL.prescribedFourierOffCenterVariance_le}
\label{FABL.prescribedFourierOffCenterVariance_le}
\leanok
\difficulty{6}
Let $n$ be a positive integer, and let $\mathcal F$ be a nonempty family of subsets of
$\{1,\ldots,n\}$. If $0<\varepsilon<1$ and $|\mathcal F|\ge (1-\varepsilon^2/(6n))2^n$,
then $2^n/|\mathcal F|-1\le \varepsilon^2/(5n)$.
\end{theorem}

\begin{theorem}[Exponential bound by an inverse power]
\lean{FABL.two_mul_exp_neg_five_halves_mul_le_four_pow_inv}
\label{FABL.two_mul_exp_neg_five_halves_mul_le_four_pow_inv}
\leanok
\difficulty{6}
For every positive natural number $n$,
\[2\exp\!\left(-\frac{5n}{2}\right)\leq\frac{1}{4^n}.\]
\end{theorem}

\begin{theorem}[Tail bound for the off-center Fourier sum]
\lean{FABL.measure_prescribedFourierOffCenterSum_abs_ge_le}
\label{FABL.measure_prescribedFourierOffCenterSum_abs_ge_le}
\leanok
\uses{FABL.BooleanFunction, FABL.measure_abs_weightedBooleanSum_ge_le, FABL.prescribedFourierKernel, FABL.prescribedFourierOffCenterSum, FABL.prescribedFourierOffCenterSum_eq_weightedSum, FABL.prescribedFourierOffCenterVariance_le, FABL.signValue, FABL.sum_sq_prescribedFourierOffCenterWeight, FABL.two_mul_exp_neg_five_halves_mul_le_four_pow_inv, FABL.uniformPMF}
\difficulty{5}
Let $n$ be a natural number, let $\mathcal F$ be a nonempty family of subsets of the $n$
coordinate indices, and let $0<\varepsilon<1$. Suppose that
$\left(1-\frac{\varepsilon^2}{6n}\right)2^n\le |\mathcal F|$, and fix a point $x$ of the
sign cube. Choose a sign-valued Boolean function uniformly at random; then the
probability that the absolute value of its prescribed Fourier off-center sum at $x$ is
at least $\varepsilon$ is at most $4^{-n}$.
\end{theorem}

\begin{definition}[Prescribed Fourier approximation]
\lean{FABL.prescribedFourierApproximation}
\label{FABL.prescribedFourierApproximation}
\leanok
\uses{FABL.BooleanFunction, FABL.prescribedFourierKernel, FABL.signValue}
For a natural number $n$, a finite family $\mathcal F$ of subsets of $\{0,\ldots,n-1\}$,
and a Boolean function $f:\{-1,1\}^n\to\{\pm1\}$, define its prescribed Fourier
approximation as the function $x\mapsto\sum_{a\in\{-1,1\}^n} f(a)K_{\mathcal F}(a,x)$,
where $K_{\mathcal F}$ is the prescribed Fourier kernel.
\end{definition}

\begin{theorem}[Vanishing coefficients outside a prescribed family]
\lean{FABL.fourierCoeff_prescribedFourierApproximation_eq_zero}
\label{FABL.fourierCoeff_prescribedFourierApproximation_eq_zero}
\leanok
\uses{FABL.BooleanFunction, FABL.fourierCoeff, FABL.fourierCoeff_prescribedFourierKernel, FABL.monomial, FABL.prescribedFourierApproximation, FABL.prescribedFourierKernel, FABL.signValue}
\difficulty{4}
Let $n$ be a nonnegative integer, let $\\mathcal F$ be a nonempty finite family of
subsets of $\\{1,\\ldots,n\\}$, let $f:\\{-1,1\\}^n\\to\\{\\pm1\\}$, and let
$T\\subseteq\\{1,\\ldots,n\\}$ with $T\\notin\\mathcal F$. For $a,x\\in\\{-1,1\\}^n$,
set $I_a(x)=\\prod_{i=1}^n(1+a_ix_i)/2$, and define $P_a(x)=\\sum_{S\\in\\mathcal
F}\\widehat I_a(S)\\prod_{i\\in S}x_i$ and $K_a(x)=2^n|\\mathcal F|^{-1}P_a(x)$. The
function $A_f(x)=\\sum_{a\\in\\{-1,1\\}^n}f(a)K_a(x)$ has Fourier coefficient zero at
$T$.
\end{theorem}

\begin{theorem}[Prescribed Fourier approximation decomposition]
\lean{FABL.prescribedFourierApproximation_apply}
\label{FABL.prescribedFourierApproximation_apply}
\leanok
\uses{FABL.BooleanFunction, FABL.BooleanFunction.toReal, FABL.prescribedFourierApproximation, FABL.prescribedFourierKernel, FABL.prescribedFourierKernel_apply_self, FABL.prescribedFourierOffCenterSum, FABL.signValue}
\difficulty{3}
Let $n$ be a natural number, let $\mathcal F$ be a nonempty finite family of subsets of
$\{0,\ldots,n-1\}$, let $f$ map the $n$-dimensional sign cube to $\{−1,1\}$, and let $x$
be a point of that cube. The prescribed Fourier approximation to $f$ at $x$ is the real
value of $f$ at $x$ plus the sum of the off-center contributions at $x$ from all other
points of the cube.
\end{theorem}

\begin{definition}[Uniform Fourier approximation property]
\lean{FABL.HasPrescribedFourierUniformApproximation}
\label{FABL.HasPrescribedFourierUniformApproximation}
\leanok
\uses{FABL.BooleanFunction, FABL.BooleanFunction.toReal, FABL.fourierCoeff}
For a natural number $n$, a finite family $\mathcal F$ of subsets of $\{0,\ldots,n-1\}$,
a real number $\varepsilon$, and a Boolean function $f:\{-1,1\}^n\to\{\pm1\}$, say that
$f$ has a uniform approximation with Fourier support in $\mathcal F$ if there is a
function $q:\{-1,1\}^n\to\mathbb R$ such that its Fourier coefficient at every subset
outside $\mathcal F$ is zero and $|f(x)-q(x)|<\varepsilon$ for every $x\in\{-1,1\}^n$.
\end{definition}

\begin{theorem}[Off-center Fourier failure bound]
\lean{FABL.measure_exists_prescribedFourierOffCenterFailure_le}
\label{FABL.measure_exists_prescribedFourierOffCenterFailure_le}
\leanok
\uses{FABL.BooleanFunction, FABL.Sign, FABL.measure_prescribedFourierOffCenterSum_abs_ge_le, FABL.prescribedFourierOffCenterSum, FABL.uniformPMF}
\difficulty{4}
Let $n$ be a natural number, let $\\mathcal F$ be a nonempty family of subsets of $[n]$,
and let $0<\\varepsilon<1$ satisfy $\\left(1-\\frac{\\varepsilon^2}{6n}\\right)2^n\\le
|\\mathcal F|$. Choose uniformly at random a sign-valued function $f$ on the Boolean
cube. The probability that there exists a point $x$ of the cube for which the absolute
value of the prescribed Fourier off-center sum associated to $\\mathcal F$, $x$, and $f$
is at least $\\varepsilon$ is at most $2^{-n}$.
\end{theorem}

\begin{theorem}[Typical approximation with prescribed Fourier support]
\lean{FABL.measure_not_hasPrescribedFourierUniformApproximation_le}
\label{FABL.measure_not_hasPrescribedFourierUniformApproximation_le}
\leanok
\uses{FABL.BooleanFunction, FABL.HasPrescribedFourierUniformApproximation, FABL.fourierCoeff_prescribedFourierApproximation_eq_zero, FABL.measure_exists_prescribedFourierOffCenterFailure_le, FABL.prescribedFourierApproximation, FABL.prescribedFourierApproximation_apply, FABL.prescribedFourierOffCenterSum, FABL.uniformPMF}
\difficulty{4}
Let $n\in\mathbb{N}$ and let $\mathcal{F}$ be a nonempty family of subsets of $[n]$.
Suppose $0<\varepsilon<1$ and $\left(1-\varepsilon^2/(6n)\right)2^n\le |\mathcal{F}|$.
For a uniformly chosen function $f:\{-1,1\}^n\to\{-1,1\}$, the probability that there is
no real-valued function $q$ whose Fourier coefficients vanish outside $\mathcal{F}$ and
which satisfies $|f(x)-q(x)|<\varepsilon$ at every input $x$ is at most $2^{-n}$.
\end{theorem}

\begin{definition}[Typical polynomial threshold cutoff]
\lean{FABL.typicalPolynomialThresholdCutoff}
\label{FABL.typicalPolynomialThresholdCutoff}
\leanok
For a natural number $n$, define the typical polynomial threshold cutoff to be the floor
of $\frac n2+\frac n2\sqrt{\frac{2\log(48n)}{n}}$.
\end{definition}

\begin{theorem}[A logarithmic bound for integers at least two]
\lean{FABL.log_48_mul_le_seven_log}
\label{FABL.log_48_mul_le_seven_log}
\leanok
\difficulty{4}
For every integer $n \ge 2$, $\log(48n) \le 7\log n$.
\end{theorem}

\begin{theorem}[Typical polynomial cutoff bound]
\lean{FABL.typicalPolynomialThresholdCutoff_le}
\label{FABL.typicalPolynomialThresholdCutoff_le}
\leanok
\uses{FABL.log_48_mul_le_seven_log, FABL.typicalPolynomialThresholdCutoff}
\difficulty{5}
For an integer $n\ge 2$, define $K_n$ to be the floor of $\frac n2+\frac
n2\sqrt{\frac{2\log(48n)}{n}}$. Then $K_n\le \frac n2+2\sqrt{n\log n}$.
\end{theorem}

\begin{theorem}[Uniform measure of a finite event]
\lean{FABL.measure_uniformPMF_real_eq_ncard_div_card}
\label{FABL.measure_uniformPMF_real_eq_ncard_div_card}
\leanok
\uses{FABL.uniformPMF}
\difficulty{4}
Let $\Omega$ be a finite nonempty measurable space with measurable singletons, and let
$A\subseteq\Omega$. The uniform probability measure of $A$ is $|A|/|\Omega|$, where
$|A|$ denotes the cardinality of $A$.
\end{theorem}

\begin{theorem}[Mean of a uniform random sign]
\lean{FABL.expect_signValue_prescribedFourierSupport}
\label{FABL.expect_signValue_prescribedFourierSupport}
\leanok
\uses{FABL.Sign, FABL.signValue}
\difficulty{1}
For a uniform random sign $s\in\{-1,1\}$, its mean is $\\mathbb{E}[s]=0$.
\end{theorem}

\begin{theorem}[Upper tail of the positive-coordinate count]
\lean{FABL.measure_positiveCoordinateCount_gt_typicalCutoff_le}
\label{FABL.measure_positiveCoordinateCount_gt_typicalCutoff_le}
\leanok
\uses{FABL.expect_signValue_prescribedFourierSupport, FABL.finiteUniformEmpiricalMean, FABL.measure_finiteUniformEmpiricalMean_sub_expect_ge_le, FABL.positiveCoordinateCount, FABL.signValue, FABL.signValue_eq_neg_one_or_one, FABL.sum_signValue_eq_two_mul_positiveCoordinateCount_sub, FABL.typicalPolynomialThresholdCutoff, FABL.uniformPMF}
\difficulty{6}
Let $n$ be a positive integer, and choose a point uniformly from the $n$-dimensional
sign cube. The probability that its number of positive coordinates exceeds $\left\lfloor
\frac n2+\frac n2\sqrt{\frac{2\log(48n)}{n}}\right\rfloor$ is at most $\frac{1}{24n}$.
\end{theorem}

\begin{theorem}[Counting vectors above the cutoff]
\lean{FABL.ncard_positiveCoordinateCount_gt_typicalCutoff}
\label{FABL.ncard_positiveCoordinateCount_gt_typicalCutoff}
\leanok
\uses{FABL.positiveCoordinateCount, FABL.signCubeEquivFinset, FABL.signCubeEquivFinset_apply_card, FABL.typicalPolynomialThresholdCutoff}
\difficulty{4}
Let $n$ be a natural number, and let $c$ be the specified typical polynomial threshold
cutoff for $n$. The number of sign vectors in $\{-1,1\}^n$ having more than $c$ positive
coordinates equals the number of subsets of an $n$-element set whose cardinality is
greater than $c$.
\end{theorem}

\begin{theorem}[Partition by Fourier degree cutoff]
\lean{FABL.card_lowDegreeFourierFamily_add_ncard_high}
\label{FABL.card_lowDegreeFourierFamily_add_ncard_high}
\leanok
\uses{FABL.lowDegreeFourierFamily, FABL.typicalPolynomialThresholdCutoff}
\difficulty{3}
For each natural number $n$, let $t_n$ be the floor of $\frac n2+\frac
n2\sqrt{\frac{2\log(48n)}{n}}$. Among all subsets of an $n$-element set, the number
having size at most $t_n$ plus the number having size greater than $t_n$ equals $2^n$.
\end{theorem}

\begin{theorem}[Size at the typical threshold cutoff]
\lean{FABL.card_lowDegreeFourierFamily_typicalPolynomialThresholdCutoff}
\label{FABL.card_lowDegreeFourierFamily_typicalPolynomialThresholdCutoff}
\leanok
\uses{FABL.Sign, FABL.card_lowDegreeFourierFamily_add_ncard_high, FABL.lowDegreeFourierFamily, FABL.measure_positiveCoordinateCount_gt_typicalCutoff_le, FABL.measure_uniformPMF_real_eq_ncard_div_card, FABL.ncard_positiveCoordinateCount_gt_typicalCutoff, FABL.positiveCoordinateCount, FABL.typicalPolynomialThresholdCutoff}
\difficulty{5}
For each natural number $n$, let $t_n=\left\lfloor\frac n2+\frac
n2\sqrt{\frac{2\log(48n)}{n}}\right\rfloor$. The number of subsets of an $n$-element set
having size at most $t_n$ is at least $\left(1-\frac{(1/2)^2}{6n}\right)2^n$.
\end{theorem}

\begin{theorem}[Low-degree approximation gives a threshold representation]
\lean{FABL.isPolynomialThreshold_of_hasLowDegreeUniformApproximation}
\label{FABL.isPolynomialThreshold_of_hasLowDegreeUniformApproximation}
\leanok
\uses{FABL.BooleanFunction, FABL.BooleanFunction.toReal, FABL.HasPrescribedFourierUniformApproximation, FABL.IsPolynomialThreshold, FABL.fourierDegree_le_iff, FABL.lowDegreeFourierFamily, FABL.thresholdSign_of_neg, FABL.thresholdSign_of_nonneg, FABL.typicalPolynomialThresholdCutoff}
\difficulty{4}
Let $n$ be a nonnegative integer and let $f:\{-1,1\}^n\to\{-1,1\}$. Set $k=\left\lfloor
\frac n2+\frac n2\sqrt{\frac{2\log(48n)}{n}}\right\rfloor$. Suppose there is a
real-valued function $q$ on the cube whose Fourier coefficients vanish on every subset
of coordinates of size greater than $k$, and such that $|f(x)-q(x)|<1/2$ for every
$x\in\{-1,1\}^n$. Then $f$ is represented by the sign of a real-valued polynomial on the
cube whose Fourier degree is at most $k$.
\end{theorem}

\begin{theorem}[Typical polynomial threshold degree bound]
\lean{FABL.measure_not_isPolynomialThreshold_typicalPolynomialThresholdCutoff_le}
\label{FABL.measure_not_isPolynomialThreshold_typicalPolynomialThresholdCutoff_le}
\leanok
\uses{FABL.BooleanFunction, FABL.HasPrescribedFourierUniformApproximation, FABL.IsPolynomialThreshold, FABL.card_lowDegreeFourierFamily_typicalPolynomialThresholdCutoff, FABL.isPolynomialThreshold_of_hasLowDegreeUniformApproximation, FABL.lowDegreeFourierFamily, FABL.lowDegreeFourierFamily_nonempty, FABL.measure_not_hasPrescribedFourierUniformApproximation_le, FABL.typicalPolynomialThresholdCutoff, FABL.uniformPMF}
\difficulty{3}
For every $n\in\mathbb{N}$, choose a uniformly random function
$f:\{-1,1\}^n\to\{-1,1\}$. The probability that $f$ has no polynomial threshold
representation of Fourier degree at most $K_n$ is at most $2^{-n}$, where $K_n$ is the
floor of $\frac n2+\frac n2\sqrt{\frac{2\log(48n)}{n}}$.
\end{theorem}
